% GENERATED FILE. Edit canonical lessons, then run npm run generate:books.
% canonical-source-hash: fnv1a64:ea98e36e786ed0d1
% canonical-lessons: MR-C347-virodh-karne, MR-C347-pathimba-dene, MR-C347-samarthan-karne, MR-C347-tatastha, MR-C347-viruddha

\chapter{For and Against}
\label{ch:mr-for-and-against}
\hlchaptermodality{347}

\begin{chapteropening}
\textbf{By the end of this chapter:} \emph{I can say who is for something and who is against it.}

Complete the last lesson of chapter 347: I can say who is for something and who is against it.
\end{chapteropening}

\section[\texorpdfstring{\mr{विरोध} \mr{करणे} (virodh karṇe)}{विरोध करणे (virodh karṇe)}]{\mr{विरोध} \mr{करणे} (virodh karṇe) — to oppose}

\label{lesson:MR-C347-virodh-karne}



\begin{quote}
Before the new one: say the Marathi for to verify, then the Marathi for to reject.
\end{quote}

\subsection*{You'll want to know: \mr{विरोध} \mr{करणे}}
\textbf{\mr{विरोध} \mr{करणे}} — \emph{virodh karṇe} — \textquotedblleft{}to oppose\textquotedblright{}.

\begin{quote}
\mr{लोकांनी} \mr{नव्या} \mr{रस्त्याला} \mr{विरोध} \mr{केला}.
\end{quote}

\begin{quote}
\emph{lokānnī navyā rastyālā virodh kelā.}
\end{quote}

\begin{quote}
People opposed the new road.
\end{quote}

Past, with -\mr{ने} (he …): \mr{त्याने} \mr{विरोध} \mr{केला}, \emph{tyāne virodh kelā}.

What is opposed takes -\mr{ला}: \mr{रस्त्याला} \mr{विरोध} \mr{करणे}.

\begin{cousinweb}
\textbf{\mr{विरोध} \mr{करणे}} (\emph{virodh karṇe}) joins the tatsama \mr{विरोध}, Sanskrit \emph{virodha}, \textquotedblleft{}opposition, hindrance\textquotedblright{} (Monier-Williams), to \mr{करणे}, \textquotedblleft{}to do\textquotedblright{}.
\end{cousinweb}

\begin{grammarlens}[title={What you've built}]
One more word for giving reasons and opinions.
\end{grammarlens}

\subsection*{Your turn}
\begin{itemize}
\raggedright
  \item \emph{Say it:} \emph{virodh karṇe}
  \item \emph{Say it:} \emph{virodh karṇe}, once more
  \item \emph{Say it:} \emph{lokānnī navyā rastyālā virodh kelā.} — People opposed the new road.
  \item \emph{Recall:} say the Marathi for to verify, then the Marathi for to reject, then say \emph{virodh karṇe} again
\end{itemize}

\begin{tcolorbox}[breakable,colback=teal!4,colframe=teal!35!black,title={Before you move on}]
What does \mr{विरोध} \mr{करणे} mean? (\textquotedblleft{}To oppose\textquotedblright{}.) Say it once more.
\end{tcolorbox}
\section[\texorpdfstring{\mr{पाठिंबा} \mr{देणे} (pāṭhimbā deṇe)}{पाठिंबा देणे (pāṭhimbā deṇe)}]{\mr{पाठिंबा} \mr{देणे} (pāṭhimbā deṇe) — to support; to back}

\label{lesson:MR-C347-pathimba-dene}



\begin{quote}
Before the new one: say the Marathi for to reject, then the Marathi for to oppose.
\end{quote}

\subsection*{You'll want to know: \mr{पाठिंबा} \mr{देणे}}
\textbf{\mr{पाठिंबा} \mr{देणे}} — \emph{pāṭhimbā deṇe} — \textquotedblleft{}to support; to back\textquotedblright{}.

\begin{quote}
\mr{माझ्या} \mr{मित्रांनी} \mr{मला} \mr{पूर्ण} \mr{पाठिंबा} \mr{दिला}.
\end{quote}

\begin{quote}
\emph{mājhyā mitrānnī malā pūrṇa pāṭhimbā dilā.}
\end{quote}

\begin{quote}
My friends gave me full support.
\end{quote}

Past, with -\mr{ने} (he …): \mr{त्याने} \mr{पाठिंबा} \mr{दिला}, \emph{tyāne pāṭhimbā dilā}.

The person backed takes -\mr{ला}: \mr{मला} \mr{पाठिंबा} \mr{दिला}.

\begin{cousinweb}
\textbf{\mr{पाठिंबा} \mr{देणे}} (\emph{pāṭhimbā deṇe}) joins \mr{पाठिंबा}, \textquotedblleft{}backing\textquotedblright{}, a Marathi noun whose origin is not settled, to \mr{देणे}, \textquotedblleft{}to give\textquotedblright{}.
\end{cousinweb}

\begin{grammarlens}[title={What you've built}]
One more word for giving reasons and opinions.
\end{grammarlens}

\subsection*{Your turn}
\begin{itemize}
\raggedright
  \item \emph{Say it:} \emph{pāṭhimbā deṇe}
  \item \emph{Say it:} \emph{pāṭhimbā deṇe}, once more
  \item \emph{Say it:} \emph{mājhyā mitrānnī malā pūrṇa pāṭhimbā dilā.} — My friends gave me full support.
  \item \emph{Recall:} say the Marathi for to reject, then the Marathi for to oppose, then say \emph{pāṭhimbā deṇe} again
\end{itemize}

\begin{tcolorbox}[breakable,colback=teal!4,colframe=teal!35!black,title={Before you move on}]
What does \mr{पाठिंबा} \mr{देणे} mean? (\textquotedblleft{}To support; to back\textquotedblright{}.) Say it once more.
\end{tcolorbox}
\section[\texorpdfstring{\mr{समर्थन} \mr{करणे} (samarthan karṇe)}{समर्थन करणे (samarthan karṇe)}]{\mr{समर्थन} \mr{करणे} (samarthan karṇe) — to justify; to defend (a view)}

\label{lesson:MR-C347-samarthan-karne}



\begin{quote}
Before the new one: say the Marathi for to oppose, then the Marathi for to support.
\end{quote}

\subsection*{You'll want to know: \mr{समर्थन} \mr{करणे}}
\textbf{\mr{समर्थन} \mr{करणे}} — \emph{samarthan karṇe} — \textquotedblleft{}to justify; to defend (a view)\textquotedblright{}.

\begin{quote}
\mr{तो} \mr{नेहमी} \mr{आपल्या} \mr{चुकीचं} \mr{समर्थन} \mr{करतो}.
\end{quote}

\begin{quote}
\emph{to nehmī āpalyā cukīcaṁ samarthan karto.}
\end{quote}

\begin{quote}
He always justifies his mistakes.
\end{quote}

Past, with -\mr{ने} (he …): \mr{त्याने} \mr{समर्थन} \mr{केलं}, \emph{tyāne samarthan kelaṁ}.

The view defended takes -\mr{चं}: \mr{निर्णयाचं} \mr{समर्थन} \mr{करणे}.

\begin{cousinweb}
\textbf{\mr{समर्थन} \mr{करणे}} (\emph{samarthan karṇe}) joins the tatsama \mr{समर्थन}, Sanskrit \emph{samarthana}, \textquotedblleft{}establishing, supporting an argument\textquotedblright{} (Monier-Williams), to \mr{करणे}, \textquotedblleft{}to do\textquotedblright{}.
\end{cousinweb}

\begin{grammarlens}[title={What you've built}]
One more word for giving reasons and opinions.
\end{grammarlens}

\subsection*{Your turn}
\begin{itemize}
\raggedright
  \item \emph{Say it:} \emph{samarthan karṇe}
  \item \emph{Say it:} \emph{samarthan karṇe}, once more
  \item \emph{Say it:} \emph{to nehmī āpalyā cukīcaṁ samarthan karto.} — He always justifies his mistakes.
  \item \emph{Recall:} say the Marathi for to oppose, then the Marathi for to support, then say \emph{samarthan karṇe} again
\end{itemize}

\begin{tcolorbox}[breakable,colback=teal!4,colframe=teal!35!black,title={Before you move on}]
What does \mr{समर्थन} \mr{करणे} mean? (\textquotedblleft{}To justify; to defend (a view)\textquotedblright{}.) Say it once more.
\end{tcolorbox}
\section[\texorpdfstring{\mr{तटस्थ} (taṭastha)}{तटस्थ (taṭastha)}]{\mr{तटस्थ} (taṭastha) — neutral}

\label{lesson:MR-C347-tatastha}



\begin{quote}
Before the new one: say the Marathi for to support, then the Marathi for to justify.
\end{quote}

\subsection*{You'll want to know: \mr{तटस्थ}}
\textbf{\mr{तटस्थ}} — \emph{taṭastha} — \textquotedblleft{}neutral\textquotedblright{}.

\begin{quote}
\mr{या} \mr{वादात} \mr{मी} \mr{तटस्थ} \mr{राहणार} \mr{आहे}.
\end{quote}

\begin{quote}
\emph{yā vādāt mī taṭastha rāhṇār āhe.}
\end{quote}

\begin{quote}
I am going to stay neutral in this dispute.
\end{quote}

Invariable: \mr{तो} \mr{तटस्थ} \mr{आहे}, \mr{ती} \mr{तटस्थ} \mr{आहे}.

\begin{cousinweb}
\textbf{\mr{तटस्थ}} (\emph{taṭastha}) is a tatsama, taken from Sanskrit \emph{taṭastha}, \textquotedblleft{}standing on the bank; standing aloof\textquotedblright{} (Monier-Williams). One who watches from the bank (\emph{taṭa}) takes no side.
\end{cousinweb}

\begin{grammarlens}[title={What you've built}]
One more word for giving reasons and opinions.
\end{grammarlens}

\subsection*{Your turn}
\begin{itemize}
\raggedright
  \item \emph{Say it:} \emph{taṭastha}
  \item \emph{Say it:} \emph{taṭastha}, once more
  \item \emph{Say it:} \emph{yā vādāt mī taṭastha rāhṇār āhe.} — I am going to stay neutral in this dispute.
  \item \emph{Recall:} say the Marathi for to support, then the Marathi for to justify, then say \emph{taṭastha} again
\end{itemize}

\begin{tcolorbox}[breakable,colback=teal!4,colframe=teal!35!black,title={Before you move on}]
What does \mr{तटस्थ} mean? (\textquotedblleft{}Neutral\textquotedblright{}.) Say it once more.
\end{tcolorbox}
\section[\texorpdfstring{\mr{विरुद्ध} (viruddha)}{विरुद्ध (viruddha)}]{\mr{विरुद्ध} (viruddha) — against; opposite}

\label{lesson:MR-C347-viruddha}



\begin{quote}
Before the new one: say the Marathi for to justify, then the Marathi for neutral.
\end{quote}

\subsection*{You'll want to know: \mr{विरुद्ध}}
\textbf{\mr{विरुद्ध}} — \emph{viruddha} — \textquotedblleft{}against; opposite\textquotedblright{}.

\begin{quote}
\mr{आमचं} \mr{मत} \mr{तुमच्या} \mr{मताच्या} \mr{अगदी} \mr{विरुद्ध} \mr{आहे}.
\end{quote}

\begin{quote}
\emph{āmcaṁ mat tumacyā matācyā agdī viruddha āhe.}
\end{quote}

\begin{quote}
Our view is the exact opposite of yours.
\end{quote}

After -\mr{च्या} or attached: \mr{तुमच्या} \mr{विरुद्ध}, \mr{संघाविरुद्ध}, against the team.

\begin{cousinweb}
\textbf{\mr{विरुद्ध}} (\emph{viruddha}) is a tatsama, taken from Sanskrit \emph{viruddha}, \textquotedblleft{}opposed, hindered\textquotedblright{} (CDIAL 11864).
\end{cousinweb}

\begin{grammarlens}[title={What you've built}]
One more word for giving reasons and opinions.
\end{grammarlens}

\subsection*{Your turn}
\begin{itemize}
\raggedright
  \item \emph{Say it:} \emph{viruddha}
  \item \emph{Say it:} \emph{viruddha}, once more
  \item \emph{Say it:} \emph{āmcaṁ mat tumacyā matācyā agdī viruddha āhe.} — Our view is the exact opposite of yours.
  \item \emph{Recall:} say the Marathi for to justify, then the Marathi for neutral, then say \emph{viruddha} again
\end{itemize}

\begin{tcolorbox}[breakable,colback=teal!4,colframe=teal!35!black,title={Before you move on}]
What does \mr{विरुद्ध} mean? (\textquotedblleft{}Against; opposite\textquotedblright{}.) Say it once more.
\end{tcolorbox}
